% =====================================================================
%  IEEE-STYLE CONFERENCE REPORT TEMPLATE  (self-contained)
%  MLE-AI-201 Machine Learning -- Final Project
%
%  This template MIMICS the IEEE two-column conference look using only
%  standard LaTeX packages. It does NOT depend on IEEEtran.cls, so it
%  compiles anywhere (tectonic, pdflatex, Overleaf) with nothing extra.
%
%  Build:   tectonic report.tex      (recommended)
%     or:   pdflatex report.tex ; pdflatex report.tex
% =====================================================================
\documentclass[10pt,twocolumn,a4paper]{article}

% ---- Geometry: IEEE-like margins / two narrow columns ----------------
\usepackage[a4paper,top=0.75in,bottom=1.0in,left=0.68in,right=0.68in,columnsep=0.25in]{geometry}

% ---- Times-like font (standard, ships with every TeX distro) ---------
\usepackage{mathptmx}          % Times text + math
\usepackage[T1]{fontenc}
\usepackage[utf8]{inputenc}    % allows accented / Azerbaijani chars

% ---- Common, dependency-free packages --------------------------------
\usepackage{graphicx}
\usepackage{amsmath,amssymb}
\usepackage{booktabs}
\usepackage{array}
\usepackage{titlesec}
\usepackage{caption}
\usepackage{cite}
\usepackage[hidelinks]{hyperref}
\usepackage{url}

% ---- IEEE-style section headings -------------------------------------
\renewcommand{\thesection}{\Roman{section}}
\renewcommand{\thesubsection}{\Alph{subsection}}
\titleformat{\section}[block]{\centering\normalsize\scshape}{\thesection.}{0.5em}{}
\titleformat{\subsection}[block]{\normalsize\itshape}{\thesubsection.}{0.5em}{}
\titlespacing*{\section}{0pt}{1.4ex plus 1ex minus .2ex}{1ex}
\titlespacing*{\subsection}{0pt}{1.0ex plus 1ex minus .2ex}{0.6ex}

% ---- Abstract / Index-Terms in the IEEE boldface-dash style ----------
\newcommand{\abstractkw}[1]{\noindent\textbf{\textit{#1}}}

\captionsetup{font=small,labelsep=period}
\setlength{\parindent}{1em}

% =====================================================================
\begin{document}

% ------------------------- TITLE BLOCK --------------------------------
\twocolumn[
  \begin{@twocolumnfalse}
  \begin{center}
    {\LARGE\bfseries Your Project Title Here:\\A Clear, Specific Subtitle\par}
    \vspace{1.2em}
    % Authors: Surname, Given name --- alphabetical by surname.
    {\large
      Aliyeva, Aysel\quad Mammadov, Murad\quad Rustamli, Rustam\par}
    \vspace{0.5em}
    {\itshape MLE-AI-201 Machine Learning, AI Academy --- Cohort II\par}
    \vspace{1.4em}
  \end{center}
  \begin{quote}
    \small
    \abstractkw{Abstract---}% IEEE uses an em-dash after "Abstract"
    One paragraph, roughly 150--250 words. State the problem (price
    prediction and the derived price-tier task on the bina.az dataset),
    what you built from scratch (a decision tree and an SVM), the key
    result (headline metrics vs.\ scikit-learn), and the main takeaway.
    Write it last. \textbf{End the abstract with your repository link},
    per IEEE convention:
    \href{https://github.com/your-team/your-repo}{\texttt{github.com/your-team/your-repo}}
    (tag \texttt{v1.0-final}).

    \vspace{0.6em}
    \abstractkw{Index Terms---}decision tree, support-vector machine,
    Pegasos, regression, classification, real-estate price prediction.
  \end{quote}
  \vspace{1.0em}
  \end{@twocolumnfalse}
]

% =====================================================================
\section{Introduction}
Motivate the problem and state your contributions as a short bulleted or
inline list. What is the task, why does it matter, and what did your
team do? Keep background brief; cite sources like this~\cite{pegasos}.

\section{Related Work / Background}
Classification and regression trees recursively partition the feature
space using threshold rules. At each node, the split is selected to reduce
task-specific impurity: Gini impurity or entropy for classification and
mean-squared error for regression~\cite{breiman1984}. This produces
interpretable piecewise-constant predictions, while tree depth and minimum
leaf size control the bias--variance trade-off. Our NumPy implementation
uses these criteria for the binary price-tier and continuous price tasks.

Support-vector machines separate classes with a maximum-margin boundary;
the soft-margin formulation trades margin width against hinge-loss
violations~\cite{cortes1995}. Pegasos optimizes the primal objective with
stochastic sub-gradient updates and a learning rate inversely proportional
to iteration and regularization strength~\cite{pegasos}. For non-linear
decision boundaries, random Fourier features approximate the RBF kernel
with an explicit finite-dimensional mapping~\cite{rahimi2007}. These
methods motivate our from-scratch Pegasos implementation and optional RFF
experiments.

\section{Data and Preprocessing}
We use the Kaggle bina.az sale-listing dataset, whose supplied CSV contains
100,775 records and 51 columns. The target is \texttt{price} in Azerbaijani
manat (AZN); property fields include location, latitude and longitude, category,
room count, area, floor, repair status, title deed, and mortgage status.
The cleaned dataset retains all 100,775 records: invalid or non-positive
prices and non-AZN listings are excluded by the cleaning rules, and exact
duplicates are removed, but none of these filters removed a record from the
provided copy. The data records AZN currency. Numeric area and room fields
are parsed from their source columns; floor and total building floors are
parsed from the floor field. Missing numeric features are imputed with
training-set medians, and missing categories are represented by an
\texttt{Unknown} category before one-hot encoding. We retain valid
high-price listings rather than applying an arbitrary outlier cap.

We exclude listing identifiers, URLs, free-text fields, scrape timestamps,
view counts, and currency labels from the model matrix. In particular,
\texttt{unit\_price} and \texttt{total\_price} are removed because they
are price-derived and would leak the regression target into both tasks.
Task A predicts price. For Task B, premium is label 1 only when price is
strictly greater than the median training price; ties are standard (0), and
the same training-only threshold labels validation and test records. The
fixed-seed-42 split is stratified by price tier and has 70\% training, 15\%
validation, and 15\% test records (70,542, 15,116, and 15,117 listings,
respectively). Standardization parameters are also fitted on training
features only. The training tier balance is 50.1\% standard and 49.9\%
premium.

\begin{table}[t]
  \centering
  \caption{Dataset summary after cleaning.}
  \label{tab:data}
  \begin{tabular}{@{}lr@{}}
    \toprule
    Property & Value \\
    \midrule
    Rows (after cleaning)       & 100,775 \\
    Features used               & 175 \\
    Train / Val / Test split    & 70 / 15 / 15\% \\
    Tier balance (\%)            & 50.1 / 49.9 \\
    \bottomrule
  \end{tabular}
\end{table}

\section{Method}
\subsection{Decision tree (from scratch)}
Describe your impurity criteria (Gini / entropy / MSE reduction), the
split search, recursion, leaf rule, and hyperparameters.

\subsection{Support-vector machine (from scratch)}
For the binary price-tier task, the two original class labels are mapped
deterministically to $\{-1,+1\}$ during fitting. Predictions are mapped
back to the original labels before evaluation. Our NumPy implementation
minimises the soft-margin primal objective
\begin{equation}
  J(\mathbf{w},b)=
  \frac{\lambda}{2}\lVert\mathbf{w}\rVert_2^2+
  \frac{1}{n}\sum_{i=1}^{n}
  \max\!\left(0,1-y_i(\mathbf{w}^{\top}\mathbf{x}_i+b)\right).
\end{equation}
Training uses single-sample stochastic Pegasos. At iteration $t$, one
training example is sampled using a fixed random seed and the learning
rate is
\begin{equation}
  \eta_t=\frac{1}{\lambda t}.
\end{equation}
The weight vector is always shrunk by $1-\eta_t\lambda$. If the sampled
margin $y_i(\mathbf{w}^{\top}\mathbf{x}_i+b)$ is below one, the update is
\begin{align}
  \mathbf{w} &\leftarrow
  (1-\eta_t\lambda)\mathbf{w}+\eta_t y_i\mathbf{x}_i,\\
  b &\leftarrow b+\eta_t y_i.
\end{align}
Otherwise, only the weight shrinkage is applied and $b$ is unchanged.
The bias is not included in the regularisation term. Inputs are required
to be numeric and standardised before reaching the SVM module.

The decision function is $f(\mathbf{x})=\mathbf{w}^{\top}\mathbf{x}+b$.
Samples satisfying $y_i f(\mathbf{x}_i)\leq 1$ are reported as support or
margin-violating samples, and the full geometric margin width is reported
as $2/\lVert\mathbf{w}\rVert_2$.

For a non-linear extension, we approximate the RBF kernel with random
Fourier features (RFF)~\cite{rahimi2007}:
\begin{equation}
  \phi(\mathbf{x})=
  \sqrt{\frac{2}{D}}
  \left[
  \cos(\boldsymbol{\omega}_1^{\top}\mathbf{x}+b_1),\ldots,
  \cos(\boldsymbol{\omega}_D^{\top}\mathbf{x}+b_D)
  \right].
\end{equation}
Each $\boldsymbol{\omega}_j$ is sampled from
$\mathcal{N}(\mathbf{0},2\gamma\mathbf{I})$, while each phase offset $b_j$
is sampled uniformly from $[0,2\pi]$. The mapping is fitted only on the
training data; validation and test data reuse exactly the same frequencies
and offsets. Linear Pegasos is then trained in the transformed space.
The number of components $D$ trades approximation quality against
computational cost, while $\gamma$ controls the locality and smoothness of
the approximated RBF kernel.

\section{Experimental Setup}
Metrics (RMSE/MAE/$R^2$ for regression; precision/recall/F1/AUC for
classification), hyperparameter search protocol, seeds, and the
scikit-learn baselines you compared against (with matched settings).

\paragraph{SVM protocol.}
The training split fits every candidate model, the validation split selects
hyperparameters using classification F1, and the test split is evaluated
exactly once after model selection. The linear search uses
$\lambda\in\{10^{-5},10^{-4},10^{-3},10^{-2},10^{-1}\}$; validation-score
ties prefer the larger $\lambda$ to favour stronger regularisation. Using
the selected $\lambda$, the RFF search evaluates
$\gamma\in\{0.01,0.1,1.0,10.0\}$; ties prefer the smaller $\gamma$. The full
experiment uses $D=256$ RFF components and $20{,}000$ Pegasos iterations.
All random seeds are fixed for reproducibility. The smaller
\texttt{fast=True} configuration is used only for development and testing,
not for final reported results. For each candidate, the experiment records
F1, final objective loss, weight norm, margin width, support count, and fit
time.

% TODO(data-ready): insert selected lambda, gamma, validation F1, test F1,
% AUC, timing, support count, and figure references from the real dataset run.

\section{Results}
Lead with a results table; back every claim with a number or figure.

\begin{table}[t]
  \centering
  \caption{Headline results (test set, mean $\pm$ std over folds).}
  \label{tab:results}
  \begin{tabular}{@{}lcc@{}}
    \toprule
    Model & Task A (RMSE) & Task B (F1) \\
    \midrule
    Ours: decision tree     & --- & --- \\
    Ours: SVM (Pegasos)     & n/a & --- \\
    scikit-learn baseline   & --- & --- \\
    \bottomrule
  \end{tabular}
\end{table}

\begin{figure}[t]
  \centering
  % \includegraphics[width=\linewidth]{figures/your_plot.pdf}
  \fbox{\parbox[c][3.2cm][c]{0.9\linewidth}{\centering\small
    Replace with a figure (e.g.\ validation curve, confusion matrix,
    or feature importances). Save plots to \texttt{figures/}.}}
  \caption{One clear figure beats ten decorative ones.}
  \label{fig:main}
\end{figure}

\section{Discussion and Limitations}
Where and why each model fails (luxury outliers, sparse districts,
scaling effects), the bias--variance story, and honest limitations.

\section{Conclusion}
One short paragraph: what you found and what you would do next.

\section*{Tools and Acknowledgements}
Disclose any AI-assistant use and what it was used for, plus data/code
you built on. (Required.)

% =====================================================================
\begin{thebibliography}{9}
\bibitem{breiman1984}
L.~Breiman, J.~H.~Friedman, R.~A.~Olshen, and C.~J.~Stone,
\emph{Classification and Regression Trees}.
Wadsworth, 1984.

\bibitem{cortes1995}
C.~Cortes and V.~Vapnik,
``Support-vector networks,''
\emph{Machine Learning}, vol.~20, pp.~273--297, 1995.

\bibitem{pegasos}
S.~Shalev-Shwartz, Y.~Singer, N.~Srebro, and A.~Cotter,
``Pegasos: primal estimated sub-gradient solver for SVM,''
\emph{Mathematical Programming}, vol.~127, no.~1, pp.~3--30, 2011.

\bibitem{rahimi2007}
A.~Rahimi and B.~Recht,
``Random features for large-scale kernel machines,''
in \emph{Advances in Neural Information Processing Systems 20},
pp.~1177--1184, 2007.
\end{thebibliography}

\end{document}
